% ===== PRÉAMBULE CANONIQUE — à coller VERBATIM en tête de CHAQUE .tex =====
\documentclass[11pt,a4paper]{article}
\usepackage[utf8]{inputenc}\usepackage[T1]{fontenc}\usepackage{lmodern}
\usepackage[french]{babel}
\usepackage{amsmath,amssymb,mathtools}\usepackage{icomma}
\usepackage{siunitx}\sisetup{locale=FR}  % \qty{}{}, \si{}, \ang{}
\usepackage{xcolor}\usepackage{colortbl}
\usepackage{tikz}\usepackage{pgfplots}\pgfplotsset{compat=1.18}
\usepackage[siunitx,europeanresistors]{circuitikz} % circuits (élec.)
\usepackage[most]{tcolorbox}\usepackage{enumitem}\usepackage{array,booktabs}
\usepackage[a4paper,margin=2.2cm]{geometry}
\usetikzlibrary{arrows.meta,calc,angles,quotes,positioning,patterns,%
  backgrounds,shapes.geometric,decorations.pathmorphing,decorations.markings,intersections}
\definecolor{primary}{RGB}{222,49,99}\definecolor{primarydark}{RGB}{150,22,55}
\definecolor{defcol}{RGB}{0,114,178}\definecolor{methcol}{RGB}{52,92,148}
\definecolor{excol}{RGB}{0,135,90}\definecolor{attncol}{RGB}{200,40,55}
\tcbset{boxbase/.style={enhanced,breakable,boxrule=0.9pt,arc=2.5pt,
  left=3mm,right=3mm,top=2.3mm,bottom=2.3mm,fonttitle=\bfseries,coltitle=white}}
% --- boîtes TUTORIEL (noms EXACTS, ne pas renommer) ---
\newtcolorbox{cle}[1][]{boxbase,colback=defcol!6,colframe=defcol,
  title={\textsc{À retenir}\ifx\relax#1\relax\relax\else\textmd{ : \itshape #1}\fi}}
\newtcolorbox{methode}[1][]{boxbase,colback=methcol!7,colframe=methcol,sharp corners,
  title={\textsc{Méthode}\ifx\relax#1\relax\relax\else\textmd{ --- \itshape #1}\fi}}
\newtcolorbox{exemplebox}[1][]{boxbase,colback=excol!5,colframe=excol,
  title={\textsc{Exemple}\ifx\relax#1\relax\relax\else\textmd{ : \itshape #1}\fi}}
\newtcolorbox{piege}[1][]{boxbase,colback=attncol!6,colframe=attncol,
  title={\textsc{Piège}\ifx\relax#1\relax\relax\else\textmd{ : \itshape #1}\fi}}
\newtcolorbox{remarque}[1][]{boxbase,colback=black!4,colframe=black!45,
  title={\textsc{Remarque}\ifx\relax#1\relax\relax\else\textmd{ : \itshape #1}\fi}}
% --- boîtes EXERCICE / CORRIGÉ (noms EXACTS, auto-comptées) ---
\newtcolorbox[auto counter]{exo}[1][]{boxbase,colback=primary!4,colframe=primary,
  title={\textsc{Exercice}~\thetcbcounter\ifx\relax#1\relax\relax\else\textmd{ --- \itshape #1}\fi}}
\newtcolorbox[auto counter]{corrige}[1][]{boxbase,colback=excol!4,colframe=excol,
  title={\textsc{Corrigé}~\thetcbcounter\ifx\relax#1\relax\relax\else\textmd{ --- \itshape #1}\fi}}
\newcommand{\vv}[1]{\vec{#1}}\newcommand{\ds}{\displaystyle}
\newcommand{\R}{\mathbb{R}}\newcommand{\N}{\mathbb{N}}\newcommand{\Z}{\mathbb{Z}}
% (un \usepackage{fancyhdr}+en-tête est possible ; il est ignoré côté web)
% =========================================================================
\begin{document}
% THEME_TITLE: Loi de Coulomb et superposition des forces
\sisetup{per-mode=symbol}

\begin{center}\large\bfseries\color{primarydark}
Thème 2 — Loi de Coulomb et superposition des forces
\end{center}

Deux charges ponctuelles exercent l'une sur l'autre une force dont \emph{Coulomb} a établi
l'intensité. Tout le thème consiste à calculer cette force, à en trouver le sens, et à
\textbf{additionner vectoriellement} plusieurs forces.

\begin{cle}[loi de Coulomb]
Deux charges ponctuelles $q_1$ et $q_2$ séparées d'une distance $d$ se repoussent (mêmes signes) ou
s'attirent (signes contraires) le long de la droite qui les joint, avec l'intensité
\[ \boxed{\,F = K\,\dfrac{|q_1|\,|q_2|}{d^{2}}\,},\qquad K \approx \num{9e9}\ \si{\newton\metre\squared\per\coulomb\squared}. \]
Les deux forces $\vv{F}_{1/2}$ et $\vv{F}_{2/1}$ ont \emph{même intensité} et sont de \emph{sens
opposés} (action--réaction), quelles que soient les valeurs de $q_1$ et $q_2$.
\end{cle}

\begin{center}
\begin{tikzpicture}[scale=1,>=Stealth,font=\footnotesize,
    ch/.style={circle,minimum size=7mm,inner sep=0pt,text=white,font=\bfseries\small}]
  % attraction
  \node[font=\bfseries,primarydark] at (-3.3,1.4) {signes opposés : attraction};
  \node[ch,fill=attncol] (a) at (-4.7,0.4) {$+$};
  \node[ch,fill=defcol] (b) at (-1.9,0.4) {$-$};
  \draw[->,line width=1.2pt,attncol] (-4.1,0.4) -- (-3.1,0.4);
  \draw[->,line width=1.2pt,defcol] (-2.5,0.4) -- (-3.5,0.4);
  \draw[<->,black!55,dashed] (-4.7,-0.35) -- (-1.9,-0.35) node[midway,below]{$d$};
  % repulsion
  \node[font=\bfseries,primarydark] at (3.3,1.4) {mêmes signes : répulsion};
  \node[ch,fill=attncol] (c) at (1.9,0.4) {$+$};
  \node[ch,fill=attncol] (e) at (4.7,0.4) {$+$};
  \draw[->,line width=1.2pt,attncol] (2.5,0.4) -- (1.5,0.4);
  \draw[->,line width=1.2pt,attncol] (4.1,0.4) -- (5.1,0.4);
  \draw[<->,black!55,dashed] (1.9,-0.35) -- (4.7,-0.35) node[midway,below]{$d$};
\end{tikzpicture}
\end{center}

\begin{cle}[trois réflexes de proportionnalité]
\begin{itemize}[leftmargin=1.5em,topsep=1pt,itemsep=1pt]
  \item $F \propto |q_1|\,|q_2|$ : doubler une charge \emph{double} $F$ ; doubler les deux la
        \emph{quadruple}.
  \item $F \propto 1/d^2$ : doubler $d$ \emph{divise} $F$ par $4$ ; tripler $d$ la divise par $9$.
  \item Diviser \emph{charges et distance par le même facteur} laisse $F$ \textbf{inchangée}
        (numérateur et dénominateur divisés pareil).
\end{itemize}
\end{cle}

\begin{methode}[superposition : additionner les forces]
La force totale sur une charge est la \textbf{somme vectorielle} des forces exercées par chacune des
autres, comme si elles étaient seules :
$\vv{F}_{\text{tot}} = \sum_i \vv{F}_{i}$.
\begin{itemize}[leftmargin=1.5em,topsep=1pt,itemsep=1pt]
  \item \textbf{Forces alignées} : on \emph{ajoute} (même sens) ou on \emph{soustrait} (sens
        opposés) les intensités.
  \item \textbf{Forces perpendiculaires} : $F_{\text{tot}} = \sqrt{F_1^{2}+F_2^{2}}$ (Pythagore),
        direction $\tan\theta = F_2/F_1$.
  \item \textbf{Symétrie} (triangle, carré, milieu) : la résultante est portée par l'axe de
        symétrie ; on projette sur cet axe.
\end{itemize}
\end{methode}

\begin{methode}[point où la force résultante s'annule (charges alignées)]
On cherche le point $M$ où les deux forces se compensent. On l'égalise :
$K\,\dfrac{|q_1|}{x^2} = K\,\dfrac{|q_2|}{(L-x)^2}$, puis on prend la \textbf{racine carrée} pour
obtenir une équation \emph{linéaire} : $\dfrac{L-x}{x} = \sqrt{\dfrac{|q_2|}{|q_1|}}$. Le point est
toujours plus près de la \emph{petite} charge.
\end{methode}

\begin{exemplebox}[application numérique directe]
$q_1 = \SI{+4}{\micro\coulomb}$, $q_2 = \SI{-1}{\micro\coulomb}$, $d = \SI{30}{\centi\metre}$ :
\[ F = \num{9e9}\times\frac{(\num{4e-6})(\num{1e-6})}{(0{,}30)^2}
     = \num{9e9}\times\frac{\num{4e-12}}{0{,}09} \approx \SI{0.40}{\newton}\ \text{(attractive)}. \]
\end{exemplebox}

\begin{piege}[les fautes qui coûtent des points]
\begin{itemize}[leftmargin=1.5em,topsep=1pt,itemsep=1pt]
  \item Oublier les \textbf{valeurs absolues} : une intensité est toujours positive ; le sens se
        décide à part.
  \item Oublier de convertir : $\si{\micro\coulomb}\to\si{\coulomb}$, $\si{\centi\metre}\to\si{\metre}$
        \emph{avant} de calculer.
  \item Additionner les intensités \emph{sans tenir compte des sens} : une somme de forces est
        \textbf{vectorielle}.
\end{itemize}
\end{piege}

\end{document}
